% META: {"id": "TNSI-HISTOIRE-INFORMATIQUE-RE-C02", "chapitre": "TNSI-HISTOIRE-INFORMATIQUE", "type_objet": "remediation", "capacites": ["T-HIST-01B"], "capacites_codes": ["C2"], "mode_creation": "ex_nihilo", "sources_inspiration": [], "status": "generated"}

%---------------------------------------------------------------
% FICHE DE REMEDIATION — C02 : croire que le materiel se specialise
%---------------------------------------------------------------

\begin{exercice}{TNSI-HIST-RE-C02-EX1}{1}{12}
Un élève écrit : « avant, une machine faisait une seule chose ; aujourd'hui les
machines font tout, donc le matériel est devenu plus spécialisé ». La conclusion
ne suit pas.
\begin{enumerate}
  \item Classer les cinq éléments suivants en \emph{matériel} ou
    \emph{logiciel} : le tableau de connexions d'un calculateur des années
    quarante, un système d'exploitation, un processeur, un navigateur Web, une
    machine virtuelle.
  \item Sur un calculateur des années quarante, changer de calcul demandait de
    recâbler la machine. Sur un téléphone actuel, changer d'usage demande
    d'installer une application. Dans lequel des deux cas le matériel est-il le
    plus \emph{général} ?
  \item Corriger la phrase de l'élève en une seule phrase, en distinguant ce qui
    s'est spécialisé de ce qui s'est généralisé.
\end{enumerate}
\end{exercice}

% BEGIN-VERIFY
% elements = {
%     "tableau de connexions": "materiel",
%     "systeme d'exploitation": "logiciel",
%     "processeur": "materiel",
%     "navigateur Web": "logiciel",
%     "machine virtuelle": "logiciel",
% }
% assert len(elements) == 5
% assert sum(1 for v in elements.values() if v == "materiel") == 2
% assert sum(1 for v in elements.values() if v == "logiciel") == 3
% assert elements["tableau de connexions"] == "materiel"
% assert elements["machine virtuelle"] == "logiciel"
% # Changer d'usage : nombre d'interventions materielles requises.
% interventions = {"calculateur 1945": 1, "telephone 2026": 0}
% assert interventions["calculateur 1945"] > interventions["telephone 2026"]
% assert interventions["telephone 2026"] == 0
% END-VERIFY

\begin{corrige}{TNSI-HIST-RE-C02-EX1}
\textbf{1.} Sont du \textbf{matériel} le tableau de connexions et le
processeur ; sont du \textbf{logiciel} le système d'exploitation, le navigateur
Web et la machine virtuelle. Deux d'un côté, trois de l'autre — et c'est déjà un
indice : la couche logicielle s'est épaissie.

\textbf{2.} C'est le téléphone dont le matériel est le plus \textbf{général}.
Changer d'usage n'y demande \emph{aucune} intervention matérielle : le processeur
exécute indifféremment un jeu, un traitement de texte ou un navigateur. Sur le
calculateur de $1945$, chaque nouveau calcul exigeait un recâblage, c'est-à-dire
une modification physique : le matériel y était l'expression même du problème
traité.

\textbf{3.} On peut écrire : « le matériel est devenu plus \emph{général}, et
c'est le logiciel qui porte désormais la spécialisation ».

Le raisonnement de l'élève confond la variété des \emph{usages} avec la
spécialisation du \emph{matériel}. Les deux évoluent en sens contraires : plus
une machine sait faire de choses, moins son matériel est adapté à l'une d'elles
en particulier. Les accélérateurs graphiques ou les circuits dédiés à
l'apprentissage automatique sont bien des spécialisations matérielles, mais ce
sont des exceptions assumées, ajoutées à côté d'un processeur généraliste, et non
un retour au câblage.
\end{corrige}
